\documentclass{minmal}
\usepackage{graphicx} % Required for inserting images
\usepackage{chemfig}
\usepackage{lewis}
\begin{document}
\textbf{\centering\large{Arbeidsoppgaver om kjemisk bining:}}\nl[1]
\liste[1] {
    \item Forklar begrepene ioniseringsenergi og elktronegativitet (her må dere ha en helt ordrett korrekt definisjon.)
    \item Forklar hvorfor elektronegativiteten stiger utover i periodene, og minker nedover i gruppene.
    \item Bestem elektronegativiteteten for disse parene av grunnstoff:
    \liste[a] {
        \item $H\quad Cl$
        \item $Na\quad Cl$
        \item $H\quad H$
        \item $C \quad H$
        \item $N \quad H$
    }
    For bindniger mellom atomer gjeler følgene regler:
    \liste[*]{
        \item Når $\Delta E<0,4$ : Da har vi elektronparbinding.
        \item når $\Delta E>0,4$ og $\Delta E<1,7$: Da har vi polar elektronparbinding.
        \item Når $\Delta E>1,7$: Da har vi ionebinding.
        
    }
    \item Hvilken binding er det mellom atom paraene over? Hint: Beregn $\Delta E$. 
    \item Tegn elektron konfigurasojon for $H \quad He \quad Na\quad S \quad O \quad Ca$
    \iten Tegn elektronprikkstruktur for $CH_4, CO_2, NH_3 og H_2O$.
    \item Er noen va molekylene i oppgven over dipol? Begrunn svaret.
}
\end{document}

